% ==== $V$ ====
\yoschapter{V}{Strong Convergence and Weak Convergence}
\label{chap:V}

In this chapter, we shall be concerned with certain basic facts pertaining to strong-, weak- and weak* convergences, including the comparison of the strong notion with the weak notion, e.g., strong- and weak measurability, and strong- and weak analyticity. We also discuss the integration of $B$-space-valued functions, that is, the theory of Bochner's integrals. The general theory of weak topologies and duality in locally convex linear topological spaces will be given in the Appendix. 


\yossection{1}{The Weak Convergence and The Weak* Convergence}
\label{sec:V.1}

\yossub{Weak Convergence}

\begin{definition}\label{def:V.1.Definition.1} $A$ sequence $\{x_{n}\}$ in a normed linear space $X$ is said to be weakly convergent if a finite $\lim\limits_{n\to \infty}f(x_n)$ exists for each $f\in X_s'$; $\{x_{n}\}$ is said to converge weakly to an element $x_{\infty}\in X$ if $\lim\limits_{n\to \infty}f(x_n) = f(x_{\infty})$ for all $f\in X_s'$. In the latter case, $x_{\infty}$ is uniquely determined, in virtue of the Hahn-Banach theorem (Corollary~\hyperref[cor:IV.6.Corollary.2]{$2$} of Theorem~\ref{thm:IV.6.Theorem.1} in Chapter~\ref{chap:IV}, \ref{sec:IV.6}); we shall write $w\text{- }\lim\limits_{n\to \infty}x_n = x_\infty$ or, in short, $x_{n}\rightarrow x_{\infty}$ weakly. $X$ is said to be sequentially weakly complete if every weakly convergent sequence of $X$ converges weakly to an element of $X$. 
\end{definition}

\begin{examplen} Let $\{x_{n}(s)\}$ be a sequence of equi-bounded continuous functions of $C[0,1]$ which is convergent to a discontinuous function $z(s)$ on $[0,1]$. Then, since $C[0,1]'$ is the space of Baire measures on $[0,1]$ of bounded total variation, we see easily that $\{x_{n}(s)\}$ gives an example of a weakly convergent sequence of $C[0,1]$ which does not converge weakly to an element of $C[0,1]$.

\end{examplen}
\begin{theorem}\label{thm:V.1.Theorem.1} \begin{enumerate}[i)]
\item $s\text{- }\lim\limits_{n \to \infty} x_n = x_\infty$ implies $w\text{- }\lim\limits_{n \to \infty} x_n = x_\infty$, but not conversely.
\item $A$ weakly convergent sequence $\{x_n\}$ is strongly bounded, and, in particular, if $w\text{- }\lim\limits_{n \to \infty} x_n = x_\infty$, then $\{\|x_n\|\}$ is bounded and $\|x_\infty\| \leqq \varliminf\limits_{n \to \infty} \|x_n\|$.
\end{enumerate}
\end{theorem}

\begin{proof}
 \begin{enumerate}[i)]
\item The first part is clear from $|f(x_n) - f(x_\infty)| \leqq \|f\| \cdot \|x_n - x_\infty\|$. The second part is proved by considering the sequence $\{x_n\}$ in the Hilbert space $(l^2)$: 

$x_{n} = \{\xi_{m}^{(n)}\}$ where $\xi_{m}^{(n)} = \delta_{nm}$ (= $1$ or $0$ according as $n = m$ or not). 

For, the value of a continuous linear functional $\in (l^2)'$ at $x = \{\xi_n\}$ is given by $\sum\limits_{n=1}^{\infty} \xi_n \bar{\eta}_n$ with some $\{\eta_n\} \in (l^2)$; consequently, $w - \lim\limits_{n \to \infty} x_n = 0$ but $\{x_n\}$ does not converge strongly to $0$ because $\|x_n\| = 1$ ( $n = 1, 2, \ldots$).
\item Consider the sequence of continuous linear functionals $X_n$ defined on the $B$-space $X_s'$ by $X_n(f) = \langle x_n, f \rangle$, and apply the resonance theorem in Chapter~\ref{chap:II}, \ref{sec:II.1}. \qedhere
\end{enumerate}
\end{proof}

\begin{theorem}[\textsc{Mazur}]\label{thm:V.1.Theorem.2} Let \wlim $_{n \to \infty} x_n = x_\infty$ in a normed linear space $X$. Then there exists, for any $\varepsilon > 0$, a convex combination $\sum\limits_{j=1}^{n} \alpha_j x_j \left( \alpha_j \geqq 0, \sum\limits_{j=1}^{n} \alpha_j = 1 \right)$ of $x_j$'s such that $\left\| x_\infty - \sum\limits_{j=1}^{n} \alpha_j x_j \right\| \leqq \varepsilon$. 
\end{theorem}

\begin{proof}
 Consider the totality $M_1$ of elements of the form $\sum\limits_{j=1}^{n} \alpha_j x_j$ with $\alpha_j \geqq 0$, $\sum\limits_{j=1}^{n} \alpha_j = 1$. We may assume that $0 \in M_1$ by replacing $x_\infty$ and $x_j$ by $(x_\infty - x_1)$ and $(x_j - x_1)$, respectively. Suppose that $\|x_\infty - u\| > \varepsilon$ for every $u \in M_1$. The set $M = \{v \in X; \|v - u\| \leqq \varepsilon/2 \text{ for some } u \in M_1\}$ is a convex neighbourhood of $0$ of $X$ and $\|x_\infty - v\| > \varepsilon/2$ for all $v \in M$. Let $p(y)$ be the Minkowski functional of $M$. Since $x_\infty = \beta^{-1} u_0$ with $p(u_0) = 1$ and $0 < \beta < 1$, we must have $p(x_\infty) = \beta^{-1} > 1$. Consider a real linear subspace $X_1 = \{x \in X; x = \gamma u_0, -\infty < \gamma < \infty\}$ and put $f_1(x) = \gamma$ for $x = \gamma u_0 \in X_1$. This real linear functional $f_1$ on $X_1$ satisfies $f(x) \leqq p(x)$ on $X_1$. Thus, by the Hahn-Banach extension theorem in Chapter~\ref{chap:IV}, \ref{sec:IV.1}, there exists a real linear extension $f$ of $f_1$ defined on the real linear space $X$ and such that $f(x) \leqq p(x)$ on $X$. $M$ being a neighbourhood of $0$, the Minkowski functional $p(x)$ is continuous in $x$. Hence $f$ is a continuous real linear functional defined on the real linear normed space $X$. We have, moreover, 
\[
\sup\limits _ {x \in M _{1}} f (x) \leqq \sup\limits _{x \in M} f (x) \leqq \sup\limits _{x \in M} p (x) = 1 < \beta^{-1} = f (\beta^{-1} u _{0}) = f (x _{\infty}). \]

Therefore, it is easy to see that $x_{\infty}$ cannot be a weak accumulation point of $M_1$, contrary to $x_{\infty} = w - \lim\limits_{n \to \infty} x_n$. 
\end{proof}

\begin{theorem}\label{thm:V.1.Theorem.3} $A$ sequence $\{x_{n}\}$ of a normed linear space $X$ converges weakly to an element $x_{\infty} \in X$ iff the following two conditions are satisfied: \begin{enumerate}[i)]
\item $\sup\limits_{n \geqq 1} \|x_{n}\| < \infty$
\item $\lim\limits_{n \to \infty} f(x_{n}) = f(x_{\infty})$ for every $f$ from any strongly dense subset $D'$ of $X_{s}'$.
\end{enumerate}
\end{theorem}

\begin{proof}
 We have only to prove the sufficiency. For any $g \in X_s'$ and $\varepsilon > 0$, there exists an $f \in D'$ such that $\| g - f \| < \varepsilon$. Thus 
\begin{align}
\left| g \left(x _{n}\right) - g \left(x _{\infty}\right) \right|& \leqq \left| g \left(x _{n}\right) - f \left(x _{n}\right) \right| + \left| f \left(x _{n}\right) - f \left(x _{\infty}\right) \right| + \left| f \left(x _{\infty}\right) - g \left(x _{\infty}\right) \right| \notag \\
&\leqq \varepsilon \| x _{n} \| + | f (x _{n}) - f (x _{\infty}) | + \varepsilon \| x _{\infty} \|, \notag
\end{align}
and hence $\lim |g(x_n) - g(x_\infty)| \leqq 2\varepsilon \sup\limits_{\infty \geqq n \geqq 1} \|x_n\|$. This proves that $\lim\limits_{n \to \infty} g(x_n) = g(x_\infty)$. 
\end{proof}

\begin{theorem}\label{thm:V.1.Theorem.4} $A$ sequence $\{x_{n}\}$ in $L^{1}(S, \mathfrak{B}, m)$ converges weakly to an element $x \in L^{1}(S, \mathfrak{B}, m)$ iff $\{\| x_{n} \|\}$ is bounded and a finite $\lim\limits_{n \to \infty} \int_{B} x_{n}(s) m(\dd s)$ exists for every $B \in \mathfrak{B}$. 
\end{theorem}

\begin{proof}
 The ``only if'' part is clear since the defining function $C_{B}(s)$ of $B \in \mathfrak{B}$ belongs to $L^{\infty}(S, \mathfrak{B}, m) = L^{1}(S, \mathfrak{B}, m)'$. 

The proof of the ``if'' part. The set function $\psi(B)=\lim\limits_{n\to\infty}\int_{B}x_{n}(s)m(d,s)$ $B\in\mathfrak{B}$, is $\sigma$-additive and $m$-absolutely continuous by the Vitali-Hahn-Saks theorem. Hence, by the differentiation theorem of \textsc{Lebesgue}-\textsc{Nikodym}, there exists an $x_{\infty} \in L^{1}(S, \mathfrak{B}, m)$ such that 
\[
\lim\limits _{n \rightarrow \infty} \int_{B} x _{n} (s) m (\dd s) = \int_{B} x _{\infty} (s) m (\dd s) \text{ for all } B \in \mathfrak {B}.
\]

Thus, for any decomposition $S = \sum\limits_{j=1}^{k} B_j$ with $B_j \in \mathfrak{B}$, we have 
\[
\lim\limits _{n \rightarrow \infty} \int_{S} x _{n} (s) y (s) m (\dd s) = \int_{S} x _{\infty} (s) y (s) m (\dd s), y (s) = \sum\limits_{j=1} ^{k} \alpha_{j} C _ {B _{j}} (s).
\]

Since such functions as $y(s)$ constitute a strongly dense subset of the space $L^{\infty}(S, \mathfrak{B}, m) = L^{1}(S, \mathfrak{B}, m)'$, we see that the ``if'' part is true by Theorem~\ref{thm:V.1.Theorem.3}. 
\end{proof}

\begin{theorem}\label{thm:V.1.Theorem.5} Let $\{x_{n}\}$ converge weakly to $x_{\infty}$ in $L^{1}(S, \mathfrak{B}, m)$. Then $\{x_{n}\}$ converges strongly to $x_{\infty}$ iff $\{x_{n}(s)\}$ converges to $x_{\infty}(s)$ in $m$-measure on every $\mathfrak{B}$-measurable set $B$ such that $m(B) < \infty$. 
\end{theorem}

\begin{remarku} $\{x_{n}(s)\}$ is said to converge to $x_{\infty}(s)$ in $m$-measure on $B$, if, for any $\varepsilon > 0$, the $m$-measure of the set $\{s \in B; |x_{n}(s) - x_{\infty}(s)| \geqq \varepsilon\}$ tends to zero as $n \to \infty$ (see the Proposition in Chapter~\ref{chap:I}, \ref{sec:I.4}). The space $(l^{1})$ is an example of $L^{1}(S, \mathfrak{B}, m)$ for which $S = \{1, 2, \ldots\}$ and $m(\{n\}) = 1$ for $n = 1, 2, \ldots$. In this case, we have $(l^{1})' = (l^{\infty})$ so that the weak convergence of $\{x_{n}\}, x_{n} = (\xi_{1}^{(n)}, \xi_{2}^{(n)}, \ldots, \xi_{k}^{(n)}, \ldots)$, to $x_{\infty} = (\xi_{1}^{(\infty)}, \xi_{2}^{(\infty)}, \ldots, \xi_{k}^{(\infty)}, \ldots)$ implies that $\lim\limits_{n \to \infty} \xi_{k}^{(n)} = \xi_{k}^{(\infty)}$ ( $k = 1, 2, \ldots$), as may be seen by taking $f \in (l^{1})'$ in such a way that $f(x) = \langle x, f \rangle = \xi_{k}$ for $x = \{\xi_{j}\} \in (l^{1})$. Thus, in the present case, $\{x_{n}\}$ converges to $x_{\infty}$ in $m$-measure on every $\mathfrak{B}$-measurable set $B$ of finite $m$-measure. In this way we obtain the 
\end{remarku}

\begin{corollaryn}[I. \textsc{Schur}] In the space $(l^1)$, if a sequence $\{x_n\}$ converges weakly to an $x_{\infty} \in (l^1)$, then \slim $x_n = x_{\infty}$.
\end{corollaryn}

\begin{proof}[Proof of Theorem~\ref{thm:V.1.Theorem.5}.]
Since the strong convergence in $L^1(S, \mathfrak{B}, m)$ implies the convergence in $m$-measure, the ``only if'' part is clear. We shall prove the ``if'' part. The sequence $\{x_n - x_\infty\}$ converges weakly to $0$, and so
\begin{equation}\label{eq:V.1.1}
\lim\limits _{n \rightarrow \infty} \int_{B} \left(x _{n} (s) - x _{\infty} (s)\right) m (\dd s) = 0 \text{ for every } B \in \mathfrak {B}.
\end{equation}

Consider the sequence of non-negative measures 
\[
\psi_{n} (B) = \int_{B} | x _{n} (s) - x _{\infty} (s) | m (\dd s), B \in \mathfrak {B}.
\]

Then we have 
\begin{equation}
\label{eq:V.1.2}
\begin{minipage}{0.92\linewidth}
\everymath{}\let\limits\nolimits
$\lim\limits_{k \to \infty} \psi_n(B_k) = 0$ uniformly in $n,$ for any decreasing sequence $\{B_k\}$ of sets $\in \mathfrak{B}$ such that $\bigcap\limits_{k=1}^{\infty} B_k = \emptyset.$
\end{minipage}
\end{equation}


If otherwise, there exists an $\varepsilon > 0$ such that for each $k$ there exists some $n_{k}$ for which $\lim\limits_{k\to \infty}n_k = \infty$ and $\psi_{n_{k}}(B_k) > \varepsilon$. Consequently, we must have 
\[
\int_ {B _{k}} \left| \operatorname{Re} \left(x _ {n _{k}} (s) - x _{\infty} (s)\right) \right| m (\dd s) > \frac {\varepsilon}{\sqrt {2}} \quad \text{ or } \quad \int_ {B _{k}} \left| \operatorname{Im} \left(x _ {n _{k}} (s) - x _{\infty} (s)\right) \right| m (\dd s) > \frac {\varepsilon}{\sqrt {2}},
\]
and so there must exist some $B_{k}^{\prime} \subseteq B_{k}$ such that 
\[
\left| \int_ {B _{k} ^{\prime}} (x _ {n _{k}} (s) - x _{\infty} (s)) m (\dd s) \right| > \varepsilon / 2 \sqrt {2} (k = 1, 2, \dots),
\]
contrary to the fact that, in virtue of \eqref{eq:V.1.1}, the $m$-absolute continuity of the sequence of measures $\varphi_{n}(B) = \int_{B} (x_{n}(s) - x_{\infty}(s)) m(\dd s)$ is uniform in $n$ (see the proof of the Vitali-Hahn-Saks Theorem in Chapter~\ref{chap:II}, \ref{sec:II.2}). 

Next let $B_0$ be any set of $\mathfrak{B}$ such that $m(B_0) < \infty$. We shall show that 
\begin{equation}\label{eq:V.1.3}
\lim\limits _{n \rightarrow \infty} \psi_{n} (B _{0}) = 0.
\end{equation}

Suppose there exist an $\varepsilon > 0$ and a subsequence $\{\psi_{n'}\}$ of $\{\psi_n\}$ such that 
\begin{equation}\label{eq:V.1.4}
\psi_ {n ^{\prime}} (B _{0}) > \varepsilon (n = 1, 2, \dots).
\end{equation}

By the hypothesis that $\{(x_n(s) - x_{\infty}(s))\}$ converges to $0$ in $m$-measure on $B_0$, there exist a subsequence $\{(x_{n''}(s) - x_{\infty}(s))\}$ of $\{(x_{n'}(s) - x_{\infty}(s))\}$ and some sets $B_n'' \subseteq B_0$ such that $m(B_n'') \leqq 2^{-n}$ and $|x_{n''}(s) - x_{\infty}(s)| < \varepsilon / m(B_0)$ on $(B_0 - B_n'')$. We put $B_k = \bigcup\limits_{n=k}^{\infty} B_n''$. Then $\{B_k\}$ is a decreasing sequence such that 
$m\left(\bigcap\limits_{k=1}^{\infty} B_k\right) \leqq \sum\limits_{n=k}^{\infty} m(B_n'') \leqq 2^{-k+1} (k=1, 2, \ldots)$ and so $m\left(\bigcap\limits_{k=1}^{\infty} B_k\right) = 0$. 

Hence, by \eqref{eq:V.1.1} and the Corollary of the Vitali-Hahn-Saks Theorem referred to above, $\lim\limits_{k\to \infty}\psi_n(B_k) = 0$ uniformly in $n$. Therefore 
\[
\psi_ {n ^{\prime \prime}} (B _{0}) \leqq \psi_ {n ^{\prime \prime}} (B _{n} ^{\prime \prime}) + \varepsilon m (B _{0}) ^{-1} \cdot m (B _{0} - B _{n} ^{\prime \prime}) \rightarrow (\leqq \varepsilon) \text{ as } n \rightarrow \infty ,
\]
contrary to \eqref{eq:V.1.4}. This proves \eqref{eq:V.1.3}. 

Now we take a sequence $\{B_k'\}$ of sets $\in \mathfrak{B}$ such that $m(B_k') < \infty$ ( $k = 1, 2, \ldots$) and $S = \bigcup\limits_{k=1}^{\infty} B_k'$. Then 
\[
\int_{S} \left| x _{n} (s) - x _{\infty} (s) \right| m (\dd s) = \int_ {\bigcup\limits_{k=1} ^{t} B _{k} ^{\prime}} + \int_ {S - \bigcup\limits_{k=1} ^{t} B _{k} ^{\prime}}
\]

By \eqref{eq:V.1.3} the first term on the right tends to zero as $n \to \infty$ for fixed $t$, and the second term on the right tends, by \eqref{eq:V.1.2}, to zero as $t \to \infty$ uniformly in $n$. Therefore we have proved that $\slim x_n = x_\infty$ in $L^1(S, \mathfrak{B}, m)$. 

A similar situation in the case of the space $\mathfrak{D}(\Omega)'$ is given by 
\end{proof}

\begin{theorem}\label{thm:V.1.Theorem.6} Let $\{T_n\}$ be a sequence of generalized functions $\in \mathfrak{D}(\Omega)'$. If $\lim\limits_{n\to \infty}T_n = T$ in the weak* topology of $\mathfrak{D}(\Omega)'$, then $\lim\limits_{n\to \infty}T_n = T$ in the strong topology of $\mathfrak{D}(\Omega)'$. 
\end{theorem}

\begin{proof}
 The strong topology of the space $\mathfrak{D}(\Omega)'$ is defined (see Definition~\ref{def:IV.7.Definition.1} in Chapter~\ref{chap:IV}, \ref{sec:IV.7}) through the family of semi-norms 

$p_{B}(T)=\sup\limits_{\varphi\in\mathfrak{B}}|T(\varphi)|$, where $B$ is any bounded set of $D$($\Omega$). 


The weak* topology of the space $\mathfrak{D}(\Omega)'$ is defined through the family of semi-norms 
$p_{F}(T)=\sup\limits_{\varphi\in\mathfrak{F}}|T(\varphi)|$, where $F$ is any finite set of $D$($\Omega$). 

Thus $\lim\limits_{n\to \infty}T_n = T$ in the weak* topology of $\mathfrak{D}(\Omega)'$ is precisely $\lim\limits_{n\to \infty}T_n = T(\mathfrak{D}(\Omega)')$ defined in Chapter~\ref{chap:II}, \ref{sec:II.3}. 

Let $\mathfrak{B}$ be any bounded set of $\mathfrak{D}(\Omega)$. Then there exists a compact subset $K$ in $\Omega$ such that $\supp(\varphi) \subseteq K$ for any $\varphi \in \mathfrak{B}$ and $\sup\limits_{x \in K, \varphi \in \mathfrak{B}} |D^j\varphi(x)|$ 
$<\infty$ for any differential operator $D^{j}$ (Theorem~\ref{thm:I.8.Theorem.1} in Chapter~\ref{chap:I}, \ref{sec:I.8}). Thus, by the Ascoli-Arzelà Theorem, $B$ is relatively compact in $\mathfrak{D}_{K}(\Omega)$. We apply the uniform boundedness theorem to the sequence $\{T_{n}-T\}$ to the effect that, for any $\varepsilon>0$, there exists a neighbourhood $U$ of $0$ of $\mathfrak{D}_{K}(\Omega)$ such that 
\[
\sup\limits _{n; \varphi \in U} | (T _{n} - T) (\varphi) | < \varepsilon .
\]

The compact subset $\mathfrak{B}$ of $\mathfrak{D}_K(\Omega)$ is covered by a finite system of sets of the form $\varphi_i + U$, where $\varphi_i \in \mathfrak{B}$ ( $i = 1, 2, \ldots, k$). Hence 
\[
\begin{aligned}
&\left| (T _{n} - T) (\varphi_{i} + u) \right|  \leqq \left| (T _{n} - T) (\varphi_{i}) \right| + \left| (T _{n} - T) (u) \right| \leqq \left| (T _{n} - T) (\varphi_{i}) \right| + \varepsilon \text{ for any } u \in U.
\end{aligned}
\]

Since $\lim\limits_{n\to \infty}(T_n - T)(\varphi_i) = 0$ for $i = 1,2,\ldots ,k$, we have 
\[
\lim\limits _{n \rightarrow \infty} (T _{n} - T) (\varphi) = 0 \text{ uniformly in } \varphi \in \mathfrak {B}.
\]

This proves our Theorem.
\end{proof}

\begin{theorem}\label{thm:V.1.Theorem.7} $A$ reflexive $B$-space $X$ is sequentially weakly complete. 
\end{theorem}

\begin{proof}
 Let a sequence $\{x_{n}\}$ of $X$ be weakly convergent. Each $x_{n}$ defines a continuous linear functional $X_{n}$ on $X_{s}^{\prime}$ by $X_{n}(x^{\prime}) = \langle x_{n}, x^{\prime} \rangle$. Since $X_{s}^{\prime}$ is a $B$-space (Theorem~\ref{thm:IV.8.Theorem.1} in Chapter~\ref{chap:IV}, \ref{sec:IV.8}), we may apply the resonance theorem. Thus a continuous linear functional on $X_{s}^{\prime}$ is defined by a finite $\lim\limits_{n \to \infty} X_{n}(x^{\prime})$ which exists by hypothesis. Since $X$ is reflexive, there exists an $x_{\infty} \in X$ such that $\langle x_{\infty}, x^{\prime} \rangle = \lim\limits_{n \to \infty} X_{n}(x^{\prime}) = \lim\limits_{n \to \infty} \langle x_{n}, x^{\prime} \rangle$, that is, $x_{\infty} = w - \lim\limits_{n \to \infty} x_{n}$. 
\end{proof}

\begin{theorem}\label{thm:V.1.Theorem.8} Let $X$ be a Hilbert space. If a sequence $\{x_{n}\}$ of $X$ converges weakly to $x_{\infty} \in X$, then $\slim\limits_{n \to \infty} x_{n} = x_{\infty}$ iff $\lim\limits_{n \to \infty} \|x_{n}\| = \|x_{\infty}\|$. 
\end{theorem}

\begin{proof}
 The ``only if'' part is clear from the continuity of the norm. The ``if'' part is clear from the equality 
\[ \begin{aligned}
&{\| x _{n} - x _{\infty} \| ^{2} = (x _{n} - x _{\infty}, x _{n} - x _{\infty})} {\qquad = \| x _{n} \| ^{2} - (x _{n}, x _{\infty}) - (x _{\infty}, x _{n}) + \| x _{\infty} \| ^{2}.}
\end{aligned}
\]

In fact, the limit, as $n \to \infty$, of the right hand side is $\| x_{\infty} \|^2 - \| x_{\infty} \|^2 - \| x_{\infty} \|^2 + \| x_{\infty} \|^2 = 0$. 
\end{proof}
\yossub{Weak* Convergence}

\begin{definition}\label{def:V.1.Definition.2} $A$ sequence $\{f_n\}$ in the dual space $X_s'$ of a normed linear space $X$ is said to be weakly* convergent if a finite $\lim\limits_{n\to \infty}f_n(x)$ exists for every $x\in X$; $\{f_n\}$ is said to converge weakly* to an element $f_\infty \in X_s'$ if $\lim\limits_{n\to \infty}f_n(x) = f_\infty (x)$ for all $x\in X$. In the latter case, we write $w^* -\lim\limits_{n\to \infty}f_n = f_\infty$ or, in short, $f_n\to f_\infty$ weakly*. 
\end{definition}

\begin{theorem}\label{thm:V.1.Theorem.9} \begin{enumerate}[i)]
\item \slim $f_n = f_\infty$ implies \wstarlim $f_n = f_\infty$, but not conversely.
\item If $X$ is a $B$-space, then a weakly* convergent sequence $\{f_n\} \subseteq X_s'$ converges weakly* to an element $f_\infty \in X_s'$ and $\|f_\infty\| \leqq \lim\limits_{n \to \infty} \|f_n\|$.
\end{enumerate}
\end{theorem}

\begin{proof}
 \begin{enumerate}[(i)]
\item The first part is clear from $|f_n(x) - f_\infty(x)| \leqq \|f_n - f_\infty\| \cdot \|x\|$. The second part is proved by the counter example given in the proof of Theorem~\ref{thm:V.1.Theorem.1}.
\item By the resonance theorem, we see that $f_\infty(x) = \lim\limits_{n \to \infty} f_n(x)$ is a continuous linear functional on $X$ and $\|f_\infty\| \leqq \lim\limits_{n \to \infty} \|f_n\|$. \qedhere
\end{enumerate}
\end{proof}

\begin{theorem}\label{thm:V.1.Theorem.10} If $X$ is a $B$-space, then a sequence $\{f_n\} \subseteq X_s'$ converges weakly* to an element $f_\infty \in X_s'$ iff \begin{enumerate}[(i)]
\item $\{\|f_n\|\}$ is bounded
\item $\lim\limits_{n \to \infty} f_n(x) = f_\infty(x)$ on a strongly dense subset of $X$.
\end{enumerate}
\end{theorem}

\begin{proof}
 The proof is similar to that of Theorem~\ref{thm:V.1.Theorem.3}. 
\end{proof}

\yossub{Strong and Weak Closure}

\begin{theorem}\label{thm:V.1.Theorem.11} Let $X$ be a locally convex linear topological space, and $M$ a closed linear subspace of $X$. Then $M$ is closed in the weak topology of $X$. 
\end{theorem}

\begin{proof}
 If otherwise, there exists a point $x_{0} \in X - M$ such that $x_{0}$ is an accumulation point of the set $M$ in the weak topology of $X$. Then, by the Corollary of Theorem~\ref{thm:IV.6.Theorem.3} in Chapter~\ref{chap:IV}, \ref{sec:IV.6}, there exists a continuous linear functional $f_{0}$ on $X$ such that $f_{0}(x_{0}) = 1$ and $f_{0}(x) = 0$ on $M$. Hence $x_{0}$ cannot be an accumulation point of the set $M$ in the weak topology of $X$. 
\end{proof}


\yossection{2}{The Local Sequential Weak Compactness of Reflexive \texorpdfstring{$B$}{B}-spaces. The Uniform Convexity}
\label{sec:V.2}

\begin{theorem}\label{thm:V.2.Theorem.1} Let $X$ be a reflexive $B$-space, and let $\{x_{n}\}$ be any sequence which is norm bounded. Then we can choose a subsequence $\{x_{n'}\}$ which converges weakly to an element of $X$. 
\end{theorem}

We will prove this Theorem under the assumption that $X$ is separable, since concrete function spaces appearing in applications of this Theorem are mostly separable. The general case of a non-separable space will be treated in the Appendix. 

\begin{lemmau} If the strong dual $X_s'$ of a normed linear space $X$ is separable, then so is $X$. 
\end{lemmau}

\begin{proof}
 Let $\{x_n'\}$ be a countable sequence which is strongly dense on the surface of the unit sphere $\{x' \in X_s'; \|x'\| = 1\}$ of $X_s'$. Choose $x_n \in X$ so that $\|x_n\| = 1$ and $|\langle x_n, x_n' \rangle| \geqq 1/2$. Let $M$ be the closed linear subspace of $X$ spanned by the sequence $\{x_n\}$. Suppose $M \neq X$ and $x_0 \in X - M$. By Corollary of Mazur's Theorem~\ref{thm:IV.6.Theorem.3} in Chapter~\ref{chap:IV}, \ref{sec:IV.6}, there exists an $x_0' \in X_s'$ such that $\|x_0'| = 1$, $\langle x_0, x_0' \rangle \neq 0$ and $\langle x, x_0' \rangle = 0$ whenever $x \in M$. Thus $\langle x_n, x_0' \rangle = 0$ ( $n = 1, 2, \ldots$), and so $1/2 \leqq |\langle x_n, x_n' \rangle| \leqq |\langle x_n, x_n' \rangle - \langle x_n, x_0' \rangle| + |\langle x_n, x_0' \rangle|$ which implies that $1/2 \leqq \|x_n\| \|x_n' - x_0'\| = \|x_n' - x_0'\|$. This is a contradiction to the fact that $\{x_n'\}$ is strongly dense on the surface of the unit sphere of $X_s'$. Thus $M = X$, and so linear combinations with rational coefficients of $\{x_n\}$ are dense in $X$. This proves our Lemma. 
\end{proof}

\begin{proof}[Proof of Theorem~\ref{thm:V.2.Theorem.1}.]
As we have remarked above, we assume that $X$ is separable and so $(X_s')_s' = X$ is separable also. By the preceding Lemma, $X_s'$ is separable. Let $\{x_n'\}$ be a countable sequence which is strongly dense in $X_s'$. Since $\{x_n\}$ is norm bounded, the sequence $\{\langle x_n, x_1' \rangle\}$ is bounded. Thus there exists a subsequence $\{x_{n_1}\}$ for which the sequence $\{\langle x_{n_1}, x_1' \rangle\}$ is convergent. Since the sequence $\{\langle x_{n_1}, x_2' \rangle\}$ is bounded, there exists a subsequence $\{x_{n_2}\}$ of $\{x_{n_2}\}$ such that $\{\langle x_{n_2}, x_2' \rangle\}$ is convergent. Proceeding in this way, we can choose a subsequence $\{x_{n_{i+1}}\}$ of the sequence $\{x_{n_{i}}\}$ such that the sequence of numbers $\{\langle x_{n_{i+1}}, x_j' \rangle\}$ converges for $j = 1, 2, \ldots, i + 1$. Hence the diagonal subsequence $\{x_{n_{n}}\}$ of the original sequence $\{x_n\}$ satisfies the condition that the sequence $\{\langle x_{n_{n}}, x_j' \rangle\}$ converges for $j = 1, 2, \ldots$. Thus, by Theorem~\hyperref[thm:V.1.Theorem.3]{$3$} in the preceding section, $\lim\limits_{n \to \infty} \langle x_{n_{n}}, x' \rangle$ exists and is finite for every $x' \in X'$. Hence, by Theorem~\hyperref[thm:V.1.Theorem.7]{$7$} of the preceding section, we see that $w - \lim\limits_{n \to \infty} x_{n_{n}}$ exists.
\end{proof}

\yossub{Milman's Theorem}

We owe to D. P. \textsc{Milman} a theorem that a $B$-space is reflexive when it is uniformly convex in the sense that, for any $\varepsilon > 0$, there exists a $\delta = \delta (\varepsilon) > 0$ such that $\|x\| \leqq 1$, $\|y\| \leqq 1$ and $\|x - y\| \geqq \varepsilon$ implies $\|x + y\| \leqq 2(1 - \delta)$. A pre-Hilbert space is uniformly convex as may be seen from the formula 
\[
\left\| x + y \right\| ^{2} + \left\| x - y \right\| ^{2} = 2 \left(\left\| x \right\| ^{2} + \left\| y \right\| ^{2}\right)
\]
valid in such a space. It is known that, for $1 < p < \infty$, the spaces $L^p$ and $(l^p)$ are uniformly convex (see J. A. \textsc{Clarkson} \cite{clarkson35}). 

\begin{theorem}[\textsc{Milman} {\cite{milman198}}]\label{thm:V.2.Theorem.2} $A$ uniformly convex $B$-space $X$ is reflexive. 
\end{theorem}

Proof (due to S. \textsc{Kakutani}). Given an $x_0'' \in (X_s')_s'$ with $\|x_0''\| = 1$. Then there exists a sequence $\{f_n\} \subseteq X_s'$ with $\|f_n\| = 1$, $x_0''(f_n) \geqq 1 - n^{-1}$ ( $n = 1, 2, \ldots$). By Theorem~\ref{thm:IV.6.Theorem.5} in Chapter~\ref{chap:IV}, \ref{sec:IV.6}, there exists, for every $n$, an $x_n \in X$ such that 
\begin{gather}\quad f _{i} (x _{n})= x _{0} ^{\prime \prime} (f _{i}) (i = 1, 2, \dots , n) \quad \text{ and } \quad \| x _{n} \| \leqq \| x _{0} ^{\prime \prime} \| + n ^{-1} = 1 + n ^{-1}. \text{ Since } \notag
\\1 - n ^{-1} \leqq x _{0} ^{\prime \prime} (f _{n}) = f _{n} (x _{n}) \leqq \| f _{n} \| \| x _{n} \| = \| x _{n} \| \leqq 1 + n ^{-1}, \notag
\end{gather}
we must have $\lim\limits_{n\to \infty}\|x_n\| = 1.$ 

If the sequence $\{x_{n}\}$ does not converge strongly, there exists an $\varepsilon > 0$ and $n_{1} < m_{1} < n_{2} < m_{2} < \cdots < n_{k} < m_{k} < \cdots$ such that $\varepsilon \leqq \|x_{n_{k}} - x_{m_{k}}\|$ ( $k = 1, 2, \ldots$). Thus, by $\lim\limits_{n \to \infty} \|x_n\| = 1$ and the uniform convexity of $X$, we obtain $\varlimsup\limits_{k \to \infty} \|x_{n_{k}} + x_{m_{k}}\| \leqq 2(1 - \delta(\varepsilon)) < 2$. But, since $n_k < m_k$, $f_{n_{k}}(x_{n_{k}}) = f_{n_{k}}(x_{m_{k}}) = x_0''(f_{n_{k}})$ and so 
\[
2 \left(1 - n _{k} ^{-1}\right) \leqq 2 x _{0} ^{\prime \prime} \left(f _ {n _{k}}\right) = f _ {n _{k}} \left(x _ {n _{k}} + x _ {m _{k}}\right) \leqq \left\| f _ {n _{k}} \right\| \cdot \left\| x _ {n _{k}} + x _ {m _{k}} \right\|.
\]

Hence, by $\| f_{n_{k}} \| = 1$, we obtain a contradiction $\varlimsup\limits_{k \to \infty} \| x_{n_{k}} + x_{m_{k}} \| \geqq 2$. 

We have thus proved the existence of $\slim\limits_{n \to \infty} x_n = x_0$, and $x_0$ satisfies 
\begin{equation}\label{eq:V.2.1}
\left\| x _{0} \right\| = 1, f _{i} (x _{0}) = x _{0} ^{\prime \prime} (f _{i}) (i = 1, 2, \dots).
\end{equation}

We show that the solution of the above equations in \eqref{eq:V.2.1} is unique. Otherwise, there exists an $\hat{x}_0 \neq x_0$ which satisfies the same equations. By the uniform convexity, $\|\hat{x}_0 + x_0\| < 2$. We also have $f_i(\hat{x}_0 + x_0) = 2x_0''(f_i)$ ( $i = 1, 2, \ldots$). Thus 
$2(1 - i^{-1}) \leqq 2x_0''(f_i) = f_i(\hat{x}_0 + x_0) \leqq \|f_i\| \|\hat{x}_0 + x_0\| = \|\hat{x}_0 + x_0\|,$ and so $\|\hat{x}_0 + x_0\| \geqq \lim\limits_{i \to \infty} 2(1 - i^{-1}) = 2$ which is a contradiction. 

Finally let $f_0$ be any point of $X_s'$. If we show that $f_0(x_0) = x_0''(f_0)$, then $(X_s')_s' \subseteq X$ and the reflexivity of $X$ is proved. To prove that $f_0(x_0) = x_0''(f_0)$, we take $f_0, f_1, \ldots, f_n, \ldots$ in place of $f_1, f_2, \ldots, f_n, \ldots$ above, and hence we obtain $\hat{x}_0 \in X$ such that 
\[
\left\| \hat {x} _{0} \right\| = 1, f _{i} (\hat {x} _{0}) = x _{0} ^{\prime \prime} (f _{i}) (i = 0, 1, \dots , n, \dots);
\]
we must have $\hat{x}_{0}=x_{0}$ by virtue of the uniqueness just proved above, and so the proof of Theorem~\ref{thm:V.2.Theorem.2} is completed. 


\yossection{3}{Dunford's Theorem and The Gelfand-Mazur Theorem}
\label{sec:V.3}

\begin{definition}\label{def:V.3.Definition.1} Let $Z$ be an open domain of the complex plane. A mapping $x(\zeta)$ defined in $Z$ with values in a $B$-space $X$ is called weakly holomorphic in $\zeta$ in the domain $Z$ if, for each $f \in X'$, the numerical function 
\[
f (x (\zeta)) = \langle x (\zeta), f \rangle
\]
of $\zeta$ is holomorphic in $Z$. 
\end{definition}

\begin{theorem}[N. \textsc{Dunford} {\cite{dunford45}}]\label{thm:V.3.Theorem.1} If $x(\zeta)$ is weakly holomorphic in $Z$, then there exists a mapping $x'(\zeta)$ defined in $Z$ with values in $X$ such that, for each $\zeta_0 \in Z$, we have 
\[
\slim _{h \to 0} h ^{-1} (x (\zeta_{0} + h) - x (\zeta_{0})) = x ^{\prime} (\zeta_{0}).
\]
\end{theorem}

In other words, the weak holomorphic property implies the strong holomorphic property. 

\begin{proof}
 Let $C$ be a rectifiable Jordan curve such that the closed bounded domain $\overline{C}$ enclosed by $C$ lies entirely in $Z$ and $\zeta_0 \in \overline{C} - C$. Let $Z_0$ be any open complex domain $\ni \zeta_0$ such that its closure lies in the interior of $\overline{C}$. Then, by Cauchy's integral representation, we have 
\[
f (x (\zeta_{0})) = \frac {1}{2 \pi i} \int_{C} \frac {f (x (\zeta))}{\zeta - \zeta_{0}} \dd \zeta . \]

Hence, if both $\zeta_0 + h$ and $\zeta_0 + g$ belong to $Z_0$, 
\[
\begin{aligned}
&\quad (h - g) ^{-1} \left\{\frac {f (x (\zeta_{0} + h)) - f (x (\zeta_{0}))}{h} - \frac {f (x (\zeta_{0} + g)) - f (x (\zeta_{0}))}{g} \right\}
\\
&= \frac {1}{2 \pi i} \int_{G} f (x (\zeta)) \left\{\frac {1}{(\zeta - \zeta_{0} - h) (\zeta - \zeta_{0} - g) (\zeta - \zeta_{0})} \right\} \dd \zeta .
\end{aligned}
\]

By the assumption, the distance between $Z_{0}$ and $C$ is positive. Hence, for fixed $f \in X'$, the absolute value of the right hand side is uniformly bounded when $\zeta_{0}, \zeta_{0} + h$ and $\zeta_{0} + g$ range over $Z_{0}$. Thus, by the resonance theorem, we have 
\[
\sup\limits _ {\zeta_{0}, \zeta_{0} + h, \zeta_{0} + g \in Z _{0}} \frac {1}{| h - g |} \left\| \left\{\frac {x (\zeta_{0} + h) - x (\zeta_{0})}{h} - \frac {x (\zeta_{0} + g) - x (\zeta_{0})}{g} \right\} \right\| < \infty .
\]

Therefore, by the completeness of the space $X$, $x(\zeta)$ is strongly differentiable at every $\zeta_0 \in Z$. 
\end{proof}

\begin{corollary}[Cauchy's integral theorem]\label{cor:V.3.Corollary.1} The strong differentiability of $x(\zeta)$ implies its strong continuity in $\zeta$. Thus we can define the curvilinear integral $\int_{C} x(\zeta) \dd\zeta$ with values in $X$. Actually we can prove that $\int_{C} x(\zeta) \dd\zeta = 0$, the zero vector of $X$. 
\end{corollary}

\begin{proof}
 We have, by the continuity and the linearity of $f \in X'$, 
\[
f \left(\int_{C} x (\zeta) \dd \zeta\right) = \int_{C} f (x (\zeta)) \dd \zeta . \]

But the right hand side is zero, because of the ordinary Cauchy integral theorem. Since $f \in X'$ was arbitrary, we must have $\int_{C} x(\zeta) \dd\zeta = 0$ by Corollary~\ref{cor:V.3.Corollary.2} of Theorem~\ref{thm:IV.6.Theorem.1} in Chapter~\ref{chap:IV}, \ref{sec:IV.6}. 
\end{proof}

From the above Corollary, we can derive other Corollaries, as in the ordinary theory of functions of a complex variable. 

\begin{corollary}[Cauchy's integral representation]\label{cor:V.3.Corollary.2} 
\[
x(\zeta_0)=\frac{1}{2\pi i}\int_{\mathcal C}\frac{x(\zeta)}{\zeta-\zeta_0}\dd\zeta
\]
for any interior point $\zeta_0$ of $\overline{\mathcal C}$.
\end{corollary}

\begin{corollary}[Taylor's expansion]\label{cor:V.3.Corollary.3} For any point $\zeta_0$ which is in the interior of the closed domain $\overline{C}$, the Taylor expansion of $x(\zeta)$ at $\zeta = \zeta_0$ converges strongly in the interior of the circle with centre at $\zeta_0$, if this circle does not extend outside of $\overline{C}$: 
\[
\begin{gathered}
x (\zeta) = \sum\limits_{n=0} ^{\infty} (n!) ^{-1} (\zeta - \zeta_{0}) ^{n} x ^{(n)} (\zeta_{0}), \text{ where } x ^{(n)} (\zeta_{0}) = \frac {n !}{2 \pi i} \int_{C} \frac {x (\zeta)}{(\zeta - \zeta_{0}) ^{n+1}} \dd \zeta .
\end{gathered}
\]
\end{corollary}

\begin{corollary}[Liouville's theorem]\label{cor:V.3.Corollary.4} If $x(\zeta)$ is (strongly) holomorphic in the whole finite plane: $|\zeta| < \infty$, and $\sup\limits_{|\zeta| < \infty} |x(\zeta)| < \infty$, then $x(\zeta)$ must reduce to a constant vector $x(0)$. 
\end{corollary}

\begin{proof}
 If we take $\left|\zeta\right|=r$ for the curve $C$, then, as $r\to\infty$, 
\[
\left\| x ^{(n)} (0) \right\| = \frac {n !}{2 \pi} \sup\limits _{| \zeta | < \infty} \left\| x (\zeta) \right\| \int_{C} \frac {| \dd \zeta |}{r ^{n+1}} \rightarrow 0 (n = 1, 2, \dots). \]

Hence the Taylor expansion of $x(\zeta)$ at $\zeta = 0$ reduces to the constant term $x(0)$ only. 
\end{proof}

We shall now apply Corollary~\ref{cor:V.3.Corollary.4} to the proof of the Gelfand-Mazur theorem. We first give 

\begin{definition}\label{def:V.3.Definition.2} $A$ commutative field $X$ over the field of complex numbers is called a normed field, if it is also a $B$-space such that the following conditions are satisfied: 

$\|e\| = 1$, where $e$ is the unit of the multiplication in $X$, 
\begin{equation}\label{eq:V.3.1}
\left\| x y \right\| \leqq \left\| x \right\| \left\| y \right\|, \text{ where } x y \text{ is the multiplication in } X.
\end{equation}
\end{definition}

\begin{theorem}[\textsc{Gelfand} {\cite{gelfand76b}}-\textsc{Mazur} {\cite{mazur191}}]\label{thm:V.3.Theorem.2} $A$ normed field $X$ is isometrically isomorphic to the complex number field. In other words, every element $x$ of $X$ is of the form $x = \xi e$ where $\xi$ is a complex number. 

\end{theorem}

\begin{proof}
 Assume the contrary, and let there exist an $x \in X$ such that $(x - \xi e) \neq 0$ for any complex number $\xi$. Since $X$ is a field, the non-zero element $(x - \xi e)$ has the inverse $(x - \xi e)^{-1} \in X$. 

We shall prove that $(x - \lambda e)^{-1}$ is (strongly) holomorphic in $\lambda$ for $|\lambda| < \infty$. We have, in fact, 
\[
\begin{gathered}
h ^{-1} ((x - (\lambda + h) e) ^{-1} - (x - \lambda e) ^{-1}) = h ^{-1} (x - (\lambda + h) e) ^{-1} \left\{e - (x - (\lambda + h) e) (x - \lambda e) ^{-1} \right\}
\\
= h ^{-1} (x - (\lambda + h) e) ^{-1} \left\{e - e + h (x - \lambda e) ^{-1} \right\} = (x - (\lambda + h) e) ^{-1} (x - \lambda e) ^{-1}.
\end{gathered}
\]

On the other hand, for sufficiently small $|h|$, the series $y^{-1}\left(e + \sum\limits_{n=1}^{\infty}(hy^{-1})^{n}\right)$, where $y = (x - \lambda e)$, converges by \eqref{eq:V.3.1}, and it represents the inverse $(y - he)^{-1} = y^{-1}(e - hy^{-1})^{-1}$, as may be seen by multiplying the series by $(y - he)$. Hence, by the strong continuity in $h$ of the series, we can prove that $(x - \lambda e)^{-1}$ is (strongly) holomorphic in $\lambda$ with the strong derivative $(x - \lambda e)^{-2}$. 

Now, if $|\lambda| \geqq 2\|x\|$, then, as above, $(x - \lambda e)^{-1} = -\lambda^{-1}(e - \lambda^{-1}x)^{-1} = -\lambda^{-1}\left(e + \sum\limits_{n=1}^{\infty}(\lambda^{-1}x)^n\right)$ and so 
\[
\left\| (x - \lambda e) ^{-1} \right\| \leqq \left| \lambda^{-1} \right|\left(1 + \sum\limits_{n=1} ^{\infty} (1 / 2) ^{n}\right)\rightarrow 0 \quad \text{ as } \quad | \lambda | \rightarrow \infty .
\]

Moreover, the function $(x - \lambda e)^{-1}$, being continuous in $\lambda$, is bounded on the compact domain of $\lambda: |\lambda| \leqq 2 \|x\|$. Hence, by Liouville's theorem, $(x - \lambda e)^{-1}$ must reduce to the constant vector $x^{-1} = (x - 0e)^{-1}$. But, since \slim $\frac{|\lambda| \to \infty}{(x - \lambda e)^{-1}} = 0$ as proved above, we have arrived at a contradiction $x^{-1} = 0$, $e = x^{-1}x = 0$. 
\end{proof}


\yossection{4}{The Weak and Strong Measurability. Pettis' Theorem}
\label{sec:V.4}

\begin{definition}\label{def:V.4.Definition.1} Let $(S, \mathfrak{B}, m)$ be a measure space, and $x(s)$ a mapping defined on $S$ with values in a $B$-space $X$. $x(s)$ is called weakly $\mathfrak{B}$-measurable if, for any $f \in X'$, the numerical function $f(x(s)) = \langle x(s), f \rangle$ of $s$ is $\mathfrak{B}$-measurable. $x(s)$ is said to be finitely-valued if it is constant $\neq 0$ on each of a finite number of disjoint $\mathfrak{B}$-measurable sets $B_j$ with $m(B_j) < \infty$ and $x(s) = 0$ on $S - \bigcup\limits_{j} B_j$. $x(s)$ is said to be strongly $\mathfrak{B}$-measurable if there exists a sequence of finitely-valued functions strongly convergent to $x(s)$ $m$-a.e. on $S$. 
\end{definition}

\begin{definition}\label{def:V.4.Definition.2} $x(s)$ is said to be separably-valued if its range $\{x(s); s \in S\}$ is separable. It is $m$-almost separably-valued if there exists a $\mathfrak{B}$-measurable set $B_0$ of $m$-measure zero such that $\{x(s); s \in S - B_0\}$ is separable. 
\end{definition}

\begin{theoremn}[B. J. \textsc{Pettis} {\cite{pettis249}}] $x(s)$ is strongly $B$-measurable iff it is weakly $B$-measurable and $m$-almost separably-valued.

\end{theoremn}
\begin{proof}
 The ``only if'' part is proved as follows. The strong $\mathfrak{B}$-measurability implies the weak $\mathfrak{B}$-measurability, because a finitely-valued function is weakly $\mathfrak{B}$-measurable, and, by the strong $\mathfrak{B}$-measurability of $x(s)$, there exists a sequence of finitely-valued functions $x_{n}(s)$ such that $s\text{- }\lim\limits_{n\to\infty}x_{n}(s)=x(s)$ except on a set $B_{0}\in\mathfrak{B}$ of $m$-measure zero. Thus the union of the ranges of $x_{n}(s)$ ( $n=1,2,\ldots$) is a countable set, and the closure of this set is separable and contains the range $\{x(s); s\in S-B_{0}\}$. 

The proof of the ``if'' part. Without losing the generality, we may assume that the range $\{x(s); s \in S\}$ is itself separable. So we may assume that the space $X$ is itself separable; otherwise, we replace $X$ by the smallest closed linear subspace containing the range of $x(s)$. We first prove that $\|x(s)\|$ is itself $B$-measurable. To this purpose, we shall make use of a lemma, to be proved later, which states that the dual space $X'$ of a separable $B$-space satisfies the condition 
\begin{equation}\label{eq:V.4.1}
\begin{minipage}[b]{0.92\linewidth}
there exists a sequence $\{f_n\} \subseteq X'$ with $\|f_n\| \leqq 1$ such that, for any $f_0 \in X'$ with $\|f_0\| \leqq 1$, we may choose a subsequence $\{f_{n'}\}$ of $\{f_n\}$ for which we have $\lim\limits_{n \to \infty} f_{n'}(x) = f_0(x)$ at every $x \in X$.
\end{minipage}
\end{equation}


Now, for any real number $a$, put 
$A = \{s; \|x(s)\| \leqq a\}$ and $A_f = \{s; |f(x(s))| \leqq a\}$, where $f \in X'$. If we can show that $A = \bigcap\limits_{j=1}^{\infty} A_{f_{j}}$, then, by the weak $B$-measurability of $x(s)$, the function $\|x(s)\|$ is $B$-measurable. It is clear that $A \subseteq_{\|f\| \leqq 1} A_f$. But, by Corollary~\hyperref[cor:IV.6.Corollary.2]{$2$} of Theorem~\ref{thm:IV.6.Theorem.1} in Chapter~\ref{chap:IV}, \ref{sec:IV.6}, there exists, for fixed $s$, an $f_0 \in X'$ with $\|f_0\| = 1$ and $f_0(x(s)) = \|x(s)\|$. Hence the reverse inclusion $A \supseteq_{\|f\| \leqq 1} A_f$ is true and so we have $A = \bigcap\limits_{\|f\| \leqq 1} A_f$. By the Lemma, we obtain $\bigcap\limits_{\|f\| \leqq 1} A_f = \bigcap\limits_{j=1}^{\infty} A_{f_{j}}$ and so $A = \bigcap\limits_{j=1}^{\infty} A_{f_{j}}$. 

Since the range $\{x(s); s \in S\}$ is separable, this range may, for any positive integer $n$, be covered by a countable number of open spheres $S_{j,n}$ ( $j = 1, 2, \ldots$) of radius $\leqq 1/n$. Let the centre of the sphere $S_{j,n}$ be $x_{j,n}$. As proved above, $\|x(s) - x_{j,n}\|$ is $\mathfrak{B}$-measurable in $s$. Hence the set $B_{j,n} = \{s \in S; x(s) \in S_{j,n}\}$ is $\mathfrak{B}$-measurable and $S = \bigcup\limits_{j=1}^{\infty} B_{j,n}$. We set 
\[
x _{n} (s) = x _{i,n} \quad \text{ if } \quad s \in B _{i,n} ^{\prime} = B _{i,n} - \bigcup\limits_{j=1} ^{i-1} B _{j,n}.
\]

Then, by $S = \sum\limits_{i=1}^{\infty} B_{i,n}'$, we have $\|x(s) - x_n(s)\| < 1/n$ for every $s \in S$.  

Since $B_{in}'$ is $\mathfrak{B}$-measurable, it is easy to see that each $x_{n}(s)$ is strongly $\mathfrak{B}$-measurable. Therefore $x(s)$, which is the strong limit of the sequence $\{x_{n}(s)\}$, is also strongly $\mathfrak{B}$-measurable. 
\end{proof}

\begin{proof}[Proof of the Lemma.]
Let a sequence $\{x_{n}\}$ be strongly dense in $X$. Consider a mapping $f \to \varphi_{n}(f) = \{f(x_{1}), f(x_{2}), \ldots, f(x_{n})\}$ of the unit sphere $S' = \{f \in X'; \|f\| \leqq 1\}$ of $X'$ into an $n$-dimensional Hilbert space $l^{2}(n)$ of vectors $(\xi_{1}, \xi_{2}, \ldots, \xi_{n})$ normed by $\left\|\left(\xi_{1}, \xi_{2}, \ldots, \xi_{n}\right)\right\| = \left(\sum\limits_{j=1}^{n} |\xi_{j}|^{2}\right)^{1/2}$. The space $l^{2}(n)$ being separable, there exists, for fixed $n$, a sequence $\{f_{n,k}\} (k = 1, 2, \ldots)$ of $S'$ such that $\{\varphi_{n}(f_{n,k}); k = 1, 2, \ldots\}$ is dense in the image $\varphi_{n}(S')$ of $S'$. 

We thus have proved that, for any $f_0 \in S'$, we can choose a subsequence $\{f_{n,m_n}\} (n = 1, 2, \ldots)$ such that $|f_{n,m_n}(x_i) - f_0(x_i)| < 1/n (i = 1, 2, \ldots, n)$. Hence $\lim\limits_{n \to \infty} f_{n,m_n}(x_i) = f_0(x_i) (i = 1, 2, \ldots)$, and so, by Theorem~\ref{thm:V.1.Theorem.10} in Chapter~\ref{chap:V}, \ref{sec:V.1}, we obtain that $\lim\limits_{n \to \infty} f_{n,m_n}(x) = f_0(x)$ for every $x \in X$. 
\end{proof}

\yossection{5}{Bochner's Integral}
\label{sec:V.5}

Let $x(s)$ be a finitely-valued function defined on a measure space $(S, \mathfrak{B}, m)$ with values in a $B$-space $X$; let $x(s)$ be equal to $x_i \neq 0$ on $B_i \in \mathfrak{B}$ ( $i = 1, 2, \ldots, n$) where $B_i$'s are disjoint and $m(B_i) < \infty$ for $i = 1, 2, \ldots, n$, and moreover, $x(s) = 0$ on $\left(S - \sum\limits_{i=1}^{n} B_i\right)$. Then we can define the $m$-integral $\int_S x(s) m(\dd s)$ of $x(s)$ over $S$ by $\sum\limits_{i=1}^{n} x_i m(B_i)$. By virtue of a limiting procedure, we can define the $m$-integral of more general functions. More precisely, we have the 

\begin{definitionu} $A$ function $x(s)$ defined on a measure space $(S, \mathfrak{B}, m)$ with values in a $B$-space $X$ is said to be Bochner $m$-integrable, if there exists a sequence of finitely-valued functions $\{x_n(s)\}$ which $s$-converges to $x(s)$ $m$-a.e. in such a way that 
\begin{equation}\label{eq:V.5.1}
\lim\limits _{n \rightarrow \infty} \int_{S} \| x (s) - x _{n} (s) \| m (\dd s) = 0.
\end{equation}
\end{definitionu}

For any set $B \in \mathfrak{B}$, the Bochner $m$-integral of $x(s)$ over $B$ is defined by 
\begin{gather}\int_{B} x (s) m (\dd s) = \slim _{n \rightarrow \infty} \int_{S} C _{B} (s) x _{n} (s) m (\dd s), \text{ where } C _{B} \text{ is the } \label{eq:V.5.2} \\\quad \text{ defining function of the set } B. \notag
\end{gather}

To justify the above definition, we have to verify that the $s$-limit on the right of \eqref{eq:V.5.2} exists and that the value of this $s$-limit is independent of the approximating sequence of functions $\{x_{n}(s)\}$. 

\begin{definitionn}[Justification of the Definition] First, $x(s)$ is strongly $\mathfrak{B}$-measurable and consequently the condition \eqref{eq:V.5.1} has a sense, since $\|x(s) - - x_n(s)\|$ is $\mathfrak{B}$-measurable as shown in the proof of Petti's Theorem. From the inequality 
\[
\begin{aligned}
&\left| \int_{B} x _{n} (s) m (\dd s) - \int_{B} x _{k} (s) m (\dd s) \right| = \left| \int_{B} (x _{n} (s) - x _{k} (s)) m (\dd s) \right|
\\
&\leqq \int_{B} \| x _{n} (s) - x _{k} (s) \| m (\dd s) \leqq \int_{S} \| x _{n} (s) - x (s) \| m (\dd s) + \int_{S} \| x (s) - x _{k} (s) \| m (\dd s)
\end{aligned}
\]

and the completeness of the space $X$, we see that $s\text{- }\lim\limits_{n\to \infty}\int_{B}x_n(s)m(\dd s)$ 
exists. It is clear also that this \slim is independent of the approximating sequence, since any two such sequences can be combined into a single approximating sequence. 
\end{definitionn}

\begin{theorem}[S. \textsc{Bochner} {\cite{bochner20}}]\label{thm:V.5.Theorem.1} $A$ strongly $\mathfrak{B}$-measurable function $x(s)$ is Bochner $m$-integrable iff $\| x(s) \|$ is $m$-integrable. 
\end{theorem}

\begin{proof}
 The ``only if'' part. We have $\left\|x(s)\right\| \leqq \left\|x_{n}(s)\right\| + \left\|x(s) - x_{n}(s)\right\|$. By the $m$-integrability of $\left\|x_{n}(s)\right\|$ and the condition \eqref{eq:V.5.1}, it is clear that $\left\|x(s)\right\|$ is $m$-integrable and 
\[
\int_{B} \| x (s) \| m (\dd s) \leqq \int_{B} \| x _{n} (s) \| m (\dd s) + \int_{B} \| x (s) - x _{n} (s) \| m (\dd s). \]

Moreover, since 
\[
\int_{B} \left| \left\| x _{n} (s) \right\| - \left\| x _{k} (s) \right\| \right| m (\dd s) \leqq \int_{B} \left\| x _{n} (s) - x _{k} (s) \right\| m (\dd s),
\]
we see from \eqref{eq:V.5.1} that $\lim\limits_{n\to \infty}\int_{B}\|x_n(s)\|m(\dd s)$ exists so that we have 
\[
\int_{B} \left\| x (s) \right\| m (\dd s) \leqq \varliminf\limits_{n \rightarrow \infty} \int_{B} \left\| x _{n} (s) \right\| m (\dd s).
\]

The ``if'' part. Let $\{x_{n}(s)\}$ be a sequence of finitely-valued functions strongly convergent to $x(s)$ $m$-a.e. Put 
\[
\begin{aligned}
&y _{n} (s) = x _{n} (s)  \text{ if } \| x _{n} (s) \| \leqq \| x (s) \| (1 + 2 ^{-1}), = 0  \text{ if } \| x _{n} (s) \| > \| x (s) \| (1 + 2 ^{-1}).
\end{aligned}
\]

Then the sequence of finitely-valued functions $\{y_n(s)\}$ satisfies $\|y_n(s)\| \leqq \|x(s)\| \cdot (1 + 2^{-1})$ and $\lim\limits_{n \to \infty} \|x(s) - y_n(s)\| = 0$ $m$-a.e. Thus, by the $m$-integrability of $\|x(s)\|$, we may apply the Lebesgue-Fatou Lemma to the functions $\|x(s) - y_n(s)\| \leqq \|x(s)\| (2 + 2^{-1})$ and obtain 
\[
\lim\limits _{n \rightarrow \infty} \int_{S} \| x (s) - y _{n} (s) \| m (\dd s) = 0,
\]
that is, $x(s)$ is Bochner $m$-integrable. 
\end{proof}
\begin{corollary}\label{cor:V.5.Corollary.1} The above proof shows that 
\[
\int_{B} \left\| x (s) \right\| m (\dd s) \geqq \left\| \int_{B} x (s) m (\dd s) \right\|,
\]
and hence $\int_{B} x(s) m(\dd s)$ is $m$-absolutely continuous in the sense that 
\[
\underset {m (B) \rightarrow 0} {\slim} \int_{B} x (s) m (\dd s) = 0.
\]
\end{corollary}

The finite additivity $\int_{\sum\limits_{j=1}^{n} B_j}^n x(s) m(\dd s) = \sum\limits_{j=1}^{n} \int_{B_{j}} x(s) m(\dd s)$ is clear and 
so, by virtue of the $\sigma$-additivity of $\int_{B} \| x(s) \| m(\dd s)$, we see that $\int_{B} x(s) m(\dd s)$ is $\sigma$-additive, i.e., 
\[
\begin{aligned}
&B = \sum\limits_{j=1} ^{\infty} B _{j} \text{ with } m (B _{j}) < \infty \text{ implies } \int_ {\sum\limits_{j=1} ^{\infty} B _{j}} x (s) m (\dd s) = \slim _{n \to \infty} \sum\limits_{j=1} ^{n} \int_ {B _{j}} x (s) m (\dd s).
\end{aligned}
\]

\begin{corollary}\label{cor:V.5.Corollary.2} Let $T$ be a bounded linear operator on a $B$-space $X$ into a $B$-space $Y$. If $x(s)$ is an $X$-valued Bochner $m$-integrable function, then $Tx(s)$ is a $Y$-valued Bochner $m$-integrable function, and 
\[
\int_{B} T x (s) m (\dd s) = T \int_{B} x (s) m (\dd s).
\]
\end{corollary}

\begin{proof}
 Let a sequence of finitely-valued functions $\{y_n(s)\}$ satisfy 
\[
\left\| y _{n} (s) \right\| \leqq \left\| x (s) \right\| (1 + n ^{-1}) \quad \text{ and } \quad \slim _{n \rightarrow \infty} y _{n} (s) = x (s) m\text{-a.e.} \]

Then, by the linearity and the continuity of $T$, we have $\int_{B} T y_n(s) m(\dd s) = T \int_{B} y_n(s) m(\dd s)$. We have, moreover, by the continuity of $T$, 
\[
\begin{aligned}&{\left\| T y _{n} (s) \right\| \leqq \left\| T \right\| \cdot \left\| y _{n} (s) \right\| \leqq \left\| T \right\| \cdot \left\| x (s) \right\| \cdot (1 + n ^{-1}) \quad \text{ and }}\\&{\qquad \slim _{n \rightarrow \infty} T y _{n} (s) = T x (s) m\text{-a.e.}}\end{aligned}
\]

Hence $Tx(s)$ is also Bochner $m$-integrable and 
\[
\begin{aligned}
&\int_{B} T x (s) m (\dd s)  = \slim\limits_{n \to \infty} \int_{B} T y _{n} (s) m (\dd s) = \slim\limits_{n \to \infty} T \int_{B} y _{n} (s) m (\dd s) = T \int_{B} x (s) m (\dd s).
\end{aligned}
\]
\end{proof}

\begin{theorem}[S. \textsc{Bochner} {\cite{bochner20}}]\label{thm:V.5.Theorem.2} Let $S$ be an $n$-dimensional euclidean space, $B$ the family of Baire sets of $S$, and $m(B)$ the Lebesgue measure of $B$. If $x(s)$ is Bochner $m$-integrable, and if $P(s_{0};\alpha)$ is the parallelopiped with centre at $s_{0}\in S$ and side length $2\alpha$, then we have the differentiation theorem 
\[
\slim _{\alpha \downarrow 0} (2 \alpha) ^{-n} \int_ {P (s _{0}; \alpha)} x (s) m (\dd s) = x (s _{0}) \text{ for } m\text{-a.e.} s _{0}.
\]
\end{theorem}

\begin{proof}
 Put 
\[
(2 \alpha) ^{-n} \int_ {P (s _{0}; \alpha)} x (s) m (\dd s) = D (x; s _{0}, \alpha). \]

If $\{x_{n}(s)\}$ is a sequence of finitely-valued functions such that $\|x_{n}(s)\| \leqq \|x(s)\| \cdot (1 + n^{-1})$ and $\slim\limits_{n\to \infty}x_n(s) = x(s)$ $m$-a.e., then 
\[
D (x; s _{0}, \alpha) - x (s _{0}) = D (x - x _{k}; s _{0}, \alpha) + D (x _{k}; s _{0}, \alpha) - x (s _{0}),
\]
and so 
\[
\begin{aligned}
\varlimsup\limits_{\alpha \downarrow 0} \| D (x; s _{0}, \alpha) - x (s _{0}) \|& \leqq \varlimsup\limits_{\alpha \downarrow 0} D (\| x - x _{k} \|; s _{0}, \alpha)
\\
& + \varlimsup\limits_{\alpha \downarrow 0} \| D (x _{k}; s _{0}, \alpha) - x _{k} (s _{0}) \| + \| x _{k} (s _{0}) - x (s _{0}) \|.
\end{aligned}
\]

The first term on the right is, by Lebesgue's theorem of differentiation of numerical functions, equal to $\|x(s_0) - x_k(s_0)\|$ $m$-a.e. The second term on the right is = $0$ $m$-a.e., since $x_k(s)$ is finitely-valued. Hence 
\[
\varlimsup\limits_{\alpha \rightarrow 0} \| D (x; s _{0}, \alpha) - x (s _{0}) \| \leqq 2 \| x _{k} (s _{0}) - x (s _{0}) \| \text{ for } m\text{-a.e.} s _{0}.
\]

Therefore, by letting $k \to \infty$, we obtain Theorem~\ref{thm:V.5.Theorem.2}. 
\end{proof}

\begin{remarku} Contrary to the case of numerical functions, a $B$-space-valued, $\sigma$-additive, $m$-absolutely continuous function need not necessarily be represented as a Bochner $m$-integral. This may be shown by a counter example. 
\end{remarku}

\begin{examplen}[$A$ Counter Example] Let $S = [0,1]$ and $\mathfrak{B}$ the family of Baire sets on $[0,1]$, and $m(B)$ the Lebesgue measure of $B \in \mathfrak{B}$. Consider the totality $m[1/3,2/3]$ of real-valued bounded functions $\xi = \xi(\theta)$ defined on the closed interval $[1/3,2/3]$ and normed by $\|\xi\| = \sup\limits_{\theta} |\xi(\theta)|$. We define an $m[1/3,2/3]$-valued function $x(s) = \xi(\theta; s)$ defined on $[0,1]$ as follows: 

the graph in $s$-$y$ plane of the real-valued function $y = y_{\theta}(s)$, which is the $\theta$-coordinate $\xi(\theta; s)$ of $x(s)$, is the polygonal line connecting the three points $(0, 0)$, $(\theta, 1)$ and $(1, 0)$ in this order. 
\end{examplen}

Then, if $s \neq s'$, we have Lipschitz' condition: 
\[
\left\| (s - s ^{\prime}) ^{-1} (x (s) - x (s ^{\prime})) \right\| = \sup\limits _{\theta} \left| (s - s ^{\prime}) ^{-1} (\xi (\theta ; s) - \xi (\theta ; s ^{\prime})) \right| \leqq 3.
\]

Thus, starting with the interval function $(x(s) - x(s'))$ taking values in $m$[$1$/$3$, $2$/$3$], we can define a $\sigma$-additive, $m$-absolutely continuous set function $x(B)$ defined for Baire set $B$ of [$0$, $1$]. 

If this function $x(B)$ is represented as a Bochner $m$-integral, then, by the preceding Theorem~\ref{thm:V.5.Theorem.2}, the function $x(s)$ must be strongly differentiable with respect to $s$ $m$-a.e. Let the corresponding strong derivative $x'(s)$ be denoted by $\eta(\theta; s)$ which takes values in $m$ [$1$/$3$, $2$/$3$]. Then, for 
every $\theta \in [1 / 3,2 / 3]$ and $m$-a.e. $s$, 
\[
\begin{aligned}
&0 = \lim\limits _{h \to 0} \| h ^{-1} (x (s + h) - x (s)) - x ^{\prime} (s) \|  \geqq \lim\limits _{h \to 0} | h ^{-1} (\xi (\theta ; s + h) - \xi (\theta ; s)) - \eta (\theta ; s) |.
\end{aligned}
\]

This proves that $\xi(\theta; s)$ must be differentiable in $s$ $m$-a.e. for all $\theta \in [1/3, 2/3]$. This is contradictory to the construction of $\xi(\theta; s)$. 

\yossub{References for Chapter V}

S. \textsc{Banach} \cite{banach9}, N. \textsc{Dunford}-J. \textsc{Schwartz} \cite{dunford44} and E. \textsc{Hille}-R. S. \textsc{Phillips} \cite{hille97}. 
